\chapter{MongoDB Indexing}

\section{What an Index Actually Does}

Without an index, MongoDB performs a \textbf{collection scan} --- checking
every document against the filter. An index is an ordered structure mapping
field values directly to documents, so MongoDB can jump straight to matches.

\begin{diagrambox}[Without an Index]
\begin{verbatim}
Query: { status: "completed" }
   Scan doc 1 -> no match
   Scan doc 2 -> no match
   Scan doc 3 -> match!
   ... (continues for every document)
\end{verbatim}
\end{diagrambox}

\begin{diagrambox}[With an Index]
\begin{verbatim}
Query: { status: "completed" }
   Look up "completed" in the status index
   Index points directly at matching documents
   Return matches (no scan of unrelated docs)
\end{verbatim}
\end{diagrambox}

\section{Unique Indexes}

A unique index speeds up lookups \emph{and} rejects duplicate values.
\texttt{unique: true} on a schema field creates it automatically.

\begin{lstlisting}[language=JavaScript]
// backend/src/models/User.js
const userSchema = new mongoose.Schema({
  email: {
    type: String,
    required: true,
    unique: true, // creates a unique index on `email`
    lowercase: true,
    trim: true,
  },
  password: { type: String, required: true },
});
\end{lstlisting}

\begin{notebox}[NOTE]
A unique index is what triggers the duplicate-email
\texttt{E11000 duplicate key error}, which your error handler can turn into
a \texttt{409 Conflict} (Chapter~25).
\end{notebox}

\section{Compound Indexes}

A compound index covers multiple fields together, valuable when a query
filters on one field and sorts on another --- the search/filter/sort pattern
from Chapter~28.

\begin{lstlisting}[language=JavaScript]
// backend/src/models/Task.js
const taskSchema = new mongoose.Schema({
  title: { type: String, required: true },
  status: { type: String, enum: ["pending", "in-progress", "completed"] },
  owner: { type: mongoose.Schema.Types.ObjectId, ref: "User" },
}, { timestamps: true });

// Speeds up: Task.find({ status: "completed" }).sort({ createdAt: -1 })
taskSchema.index({ status: 1, createdAt: -1 });

module.exports = mongoose.model("Task", taskSchema);
\end{lstlisting}

Field order matters: this index efficiently supports filtering on
\texttt{status} alone or \texttt{status} + sort by \texttt{createdAt}, but
not a sort by \texttt{createdAt} without the \texttt{status} filter.

\section{When to Add an Index}

\begin{notebox}[NOTE]
Index fields used frequently to \textbf{filter} (\texttt{status},
\texttt{owner}), \textbf{sort} (\texttt{createdAt}), or \textbf{look up
uniquely} (\texttt{email}).
\end{notebox}

\begin{warnbox}[WARNING]
Indexes aren't free --- each one slows down writes and uses extra disk and
memory, so only index fields that are actually queried or sorted on often.
\end{warnbox}
